\documentclass{article}
\newcommand{\abold}{{\bf a}}
\newcommand{\bbold}{{\bf b}}
\newcommand{\cbold}{{\bf c}}
\newcommand{\dbold}{{\bf d}}
\newcommand{\ebold}{{\bf e}}
\newcommand{\fbold}{{\bf f}}
\newcommand{\gbold}{{\bf g}}
\newcommand{\hbold}{{\bf h}}
\newcommand{\ibold}{{\bf i}}
\newcommand{\jbold}{{\bf j}}
\newcommand{\kbold}{{\bf k}}
\newcommand{\lbold}{{\bf l}}
\newcommand{\mbold}{{\bf m}}
\newcommand{\nbold}{{\bf n}}
\newcommand{\obold}{{\bf o}}
\newcommand{\pbold}{{\bf p}}
\newcommand{\qbold}{{\bf q}}
\newcommand{\rbold}{{\bf r}}
\newcommand{\sbold}{{\bf s}}
\newcommand{\tbold}{{\bf t}}
\newcommand{\ubold}{{\bf u}}
\newcommand{\vbold}{{\bf v}}
\newcommand{\wbold}{{\bf w}}
\newcommand{\xbold}{{\bf x}}
\newcommand{\ybold}{{\bf y}}
\newcommand{\zbold}{{\bf z}}


\newcommand{\half}{{\frac{1}{2}}}              
\newcommand{\iphf}{{\ibold  + \half\ebold^{d_f}}}
\newcommand{\ipf} {{\ibold  +      \ebold^{d_f}}}
\newcommand{\ipo}{{\ibold   +      \ebold^{d_d}}}
\newcommand{\imo}{{\ibold   -      \ebold^{d_d}}}
\newcommand{\lolo}{{\ibold  -      \ebold^{d_d}}}
\newcommand{\hilo}{{\ibold  +      \ebold^{d_d}}}
\newcommand{\lohi}{{\ibold  -      \ebold^{d_d}  +      \ebold^{d_f}}}
\newcommand{\hihi}{{\ibold  +      \ebold^{d_d}  +      \ebold^{d_f}}}

\begin{document}

% Title, authors and addresses

\title{Viscous Tensor Operator with spatially varying coeffcients}

\author{D. T. Graves}

\section{Introduction}
  This is a specification for the viscous tensor operator used
  in the ice sheet code.

\section{Equations}

Given a charge distribution $\rho$, The viscous tensor equation for
the velocity $v$ is given by
$$
\alpha v + \beta \nabla \cdot F(v) = \rho,
$$
where $\alpha$ and $\beta$ are constants and the flux tensor $F$ is
given by
$$
F(v) = \mu(x) (\nabla v + (\nabla v)^T) + \lambda(x) I 
$$
where $I$ is the identity matrix and $\mu$ and $\lambda$ are scalar
coefficients that vary in space.   Typically $\lambda(x) = C \mu(x)$ for
some constant $C$.

\section{Discretization}

Normal gradients at faces are computed with  centered differences.
The components of the gradient tangential to the face are computed
using an average of surrounding cell-centered gradients.   Given a
scalar field $\phi$, we compute the gradient at face ($\iphf$) and grid
spacing $h$ as
$$
G^{d_d}(\phi)_{\iphf} = (\phi_{\ipf}-\phi_{\ibold})/{h}
$$

%%  \mbox{if $d_d = d_f$},
$$
G^{d_d}(\phi)_\iphf = \frac{1}{4}(G^{d_d}_{\lolo} + G^{d_d}_{\lohi} +G^{d_d}_{\hihi}  + G^{d_d}_{\hilo})   \mbox{if $d_d \ne d_f$},
$$

where

$$
G^{d_d}(phi)_\ibold = \frac{\phi_{\ipo} - \phi_{\imo}}{2 h}
$$


\end{document}

